\chapter{Impact of the Mask Generator on Instance-Level Mapping}
\label{ch:instance_tracking}

% EJE AGNÓSTICO. El segmentador acota cantidad (recall) y calidad (precision) del mapa
% antes de cualquier razonamiento 3D. Pregunta del capítulo: ¿un segmentador más nuevo
% (SAM3) mejora ambos sin cambiar OVO? Handoff al contest (\ref{ch:contest}): cambiar de
% modelo no rompe el techo del paradigma AMG (grid + poda heurística, máscara de un tiro).

El sistema descrito en el Capítulo~\ref{ch:preliminaries} construye el mapa de instancias a
partir de las máscaras 2D que un segmentador produce cada ciertos keyframes. Esa dependencia es
total y unidireccional: la asignación de instancias toma las máscaras como dato de entrada y, una vez
asignadas las instancias, no revisa su decisión en ningún momento posterior. Esa decisión se propaga al
mapa sin posibilidad de corrección. La confianza del sistema en el segmentador es, por tanto, ciega, rígida
e incorregible, y la calidad del mapa final queda acotada por la del segmentador antes de que intervenga
ningún razonamiento 3D. Esta dependencia se manifiesta en dos ejes, la \textbf{cantidad} y la
\textbf{calidad} de las máscaras, ilustrados respectivamente en las
Figuras~\ref{fig:seg_quantity} y~\ref{fig:seg_quality}.

\begin{figure}[htbp]
  \centering
  \includegraphics[width=\textwidth]{seg_quantity}
  \caption[El eje de cantidad]{El eje de \textbf{cantidad}, con el segmentador base (SAM2).
    Arriba, las máscaras 2D; abajo, la instancia 3D resultante. Los casos \ding{192}--\ding{195}
    se comentan en el texto.}
  \label{fig:seg_quantity}
\end{figure}

La \textbf{cantidad}, entendida como el número y la granularidad de las máscaras, fija cuántas
oportunidades hay de crear o refinar instancias, y falla de dos formas.

La primera es \emph{por defecto}. Un punto 3D solo puede crear una instancia o asociarse a ella
si alguna máscara 2D lo cubre; los píxeles sin máscara no participan en la construcción semántica del
mapa. Si el segmentador no marca un objeto, este puede no llegar nunca a ser instancia semantica 3D, y su
única salida es que otro keyframe lo segmente, lo que no siempre ocurre. En la
Figura~\ref{fig:seg_quantity}, SAM2 no cubre un trozo de pared (\ding{192}) ni una parte del
sofá (\ding{193}). El desenlace depende de si esa región es segmentada más adelante: la pared,
que no vuelve a verse en ningún otro frame, queda definitivamente sin asignar, mientras que la
parte del sofá se recupera parcialmente al segmentarse en una observación posterior.

La segunda es \emph{por exceso de fragmentación}, el problema de granularidad. Un mismo objeto
puede salir troceado en varias máscaras y generar varias instancias para una única entidad. Lo
decisivo es cómo se observa por primera vez: si el objeto se ve entero, se reconstruye como una
sola instancia aunque más tarde reaparezca fragmentado; pero si la primera vez se ve a trozos,
nace dividido, y ningún mecanismo posterior une esos fragmentos, porque no llegan a solaparse
entre sí. Así ocurre en la Figura~\ref{fig:seg_quantity} con el sofá (\ding{194}) y el suelo
(\ding{195}): este último, pese a estar bien delimitado en 2D, aparece dividido en tres
instancias en 3D porque la primera vez que se segmentó estaba troceado.

\begin{figure}[htbp]
  \centering
  \includegraphics[width=\textwidth]{seg_quality}
  \caption[El eje de calidad]{El eje de \textbf{calidad}, con el segmentador base (SAM2). Arriba,
    la máscara 2D; abajo, la instancia 3D contaminada. Los casos \ding{192}--\ding{194} se
    comentan en el texto.}
  \label{fig:seg_quality}
\end{figure}

El segundo eje es la \textbf{calidad}, la precisión con que cada máscara delimita su objeto. Una
máscara que invade regiones vecinas arrastra ese error de forma permanente: los puntos mal
atribuidos quedan asignados a la instancia equivocada y ya no pueden reasociarse a la correcta,
porque pertenecen a otra. La contaminación puede ser incluso lo bastante fuerte como para que,
cuando esos objetos se segmenten bien más adelante, el voto por mayoría que decide la asignación
los siga considerando parte de la instancia intrusa.

En la Figura~\ref{fig:seg_quality}, la máscara de la pantalla de la pared (\ding{192}) se
derrama sobre dos objetos que no le pertenecen: la mesa (\ding{193}) y la silla de escritorio
(\ding{194}). El efecto sobre cada uno difiere. La mesa, aunque después se segmenta de forma
individual, hereda la contaminación inicial: buena parte de sus puntos se atribuyen a la
pantalla, como si formaran parte de ella, y ese error degrada tanto las métricas class-agnostic
como las semánticas. En la silla, la contaminación es más parcial: solo los puntos afectados en
la primera observación siguen marcados como parte de la pantalla, mientras que el resto se
reconstruye correctamente. En ambos casos, la delimitación defectuosa de la máscara inicial es
irreversible y su error se propaga al mapa final.

El flujo que va de la imagen al mapa se descompone en tres pasos: la detección de máscaras, que
en cada frame delimita las regiones correspondientes a objetos; el seguimiento, que mantiene su
identidad a lo largo del tiempo; y la asignación, que traslada las máscaras a la nube de puntos
y decide a qué instancia 3D pertenece cada uno, y que OVO resuelve mediante la asociación
temporal descrita en el Capítulo~\ref{ch:preliminaries}. Los dos ejes anteriores nacen en el
primer paso, la máscara, y por eso un mismo problema admite dos ataques complementarios. El
primero, objeto de este capítulo, es mejorar la materia prima: obtener máscaras más completas y
mejor delimitadas cambiando el segmentador. El segundo es hacer más con las máscaras
disponibles, actuando sobre la asignación para corregir aguas abajo los defectos que la
detección no evita. Este capítulo aborda el primero; el segundo se desarrolla a partir del
Capítulo~\ref{ch:contest}.

\section{Segmentation and the Automatic Mask Generator (AMG)}
\label{sec:seg_amg}

La pieza que queremos mejorar no es el segmentador entero, sino una parte muy concreta de él: el
mecanismo que barre cada keyframe en busca de máscaras 2D. Su granularidad y su exhaustividad
fijan qué objetos pueden llegar a existir en el mapa. Los segmentadores de la familia SAM se
emplean en OVO precisamente porque saben segmentar cualquier objeto, sin compromiso semántico.
Ahora bien, su uso habitual parte de un prompt que indica qué objeto seguir, y OVO no dispone de
esa información: no sabe de antemano qué objetos hay en la escena ni cuántos son. Para
descubrirlos por sí mismo, sin ningún prompt externo, se recurre al mecanismo incluido en SAM2 de
generación automática de máscaras (\emph{Automatic Mask Generator}, AMG).

El mecanismo AMG se aplica sobre cada keyframe. La imagen RGB se pasa una sola vez por el encoder
de visión de SAM2, que produce un feature map, y sobre el frame se dispone una rejilla regular de
puntos que actúan como prompts geométricos. Reutilizando ese mismo feature map, el mask decoder
propone varias máscaras por punto y, tras un filtrado interno por confianza, conserva las que
superan el umbral. El resultado es un conjunto de máscaras 2D que cubre la escena, todavía
redundante y con solapes y contenciones. Sobre él, OVO aplica una segunda poda heurística que
descarta máscaras según umbrales fijos de solape, contención y nota.

\begin{figure}[htbp]
  \centering
  \includegraphics[width=\textwidth]{amg}
  \caption[El AMG de OVO con SAM2, paso a paso]{El AMG en OVO con SAM2, sobre el frame 70 de
    \emph{office0} de Replica.}
  \label{fig:amg_pipeline}
\end{figure}

La Figura~\ref{fig:amg_pipeline} muestra el proceso sobre un keyframe de Replica: de los
$16\times16=256$ puntos de la rejilla salen 23 máscaras candidatas, que la poda deja en 19. La
mayoría de los descartes son legítimos: duplicados exactos, como las máscaras 5 y 18, o máscaras
contenidas en otra mejor, como el símbolo de reciclado (máscara 11), absorbido por la máscara del
cubo de basura que lo contiene (máscara 6). El segmentador incluso resuelve casos difíciles: la
máscara 19 reúne en una sola dos fragmentos de la alfombra que una oclusión parcial mantiene
separados en la imagen, algo que la asignación por sí sola no podría recomponer; y aunque la
máscara 17 solo captura uno de esos fragmentos, la poda conserva la 19, más completa.

Pero la criba es puramente geométrica y no distingue un solapamiento benigno de uno dañino. La
máscara 22 abarca a la vez la pizarra y la mesa y, al solaparse, elimina la máscara 4, que
aislaba correctamente la mesa: sobrevive la versión fusionada, que en 3D creará y expandirá una
única instancia donde deberían existir dos. En cualquier caso, ninguna de estas decisiones se
revisa después: la máscara que queda es la que la asignación hereda y sobre la que se construye
la instancia. Cada máscara es, así, una apuesta de un solo intento.

El objetivo es determinar si, con el mismo mecanismo pero un segmentador más reciente, y sin
alterar el resto de OVO, se puede mejorar tanto la cantidad como la calidad de las máscaras, y
por tanto la calidad del mapa final. La propuesta se centra en examinar si SAM3, el modelo más
reciente de la familia, mejora el trabajo de su predecesor. De confirmarse la mejora con solo ese
cambio, SAM3 ofrece además un incentivo: al integrar de fábrica detección y seguimiento, abre la
vía a resolver con un único mecanismo varios problemas que OVO aborda hoy por separado.

\subsection{SAM3: architecture and differences from SAM2}
\label{sec:seg_sam3}

SAM3 es la evolución de la familia SAM que incorpora comprensión semántica. Donde SAM2 usaba
como backbone de visión un modelo Hiera entrenado con masked autoencoding, cuyas features son
puramente visuales, SAM3 emplea un Perception Encoder (PE), un modelo visión-lenguaje entrenado
con una pérdida contrastiva imagen-texto. Sus features quedan así alineadas con el espacio del
lenguaje, lo que permite ir más allá de las máscaras previas y los prompts geométricos: al
añadir un text encoder y un detector basado en DETR condicionado por texto, SAM3 admite prompts
de texto y habilita la segmentación por concepto (\emph{open-vocabulary}).

Esa capacidad semántica, sin embargo, es ajena al problema que nos ocupa. Alimentar el mapa
exige detectar en cada keyframe el mayor número posible de objetos sin conocimiento previo de la
escena, una tarea que solo puede abordarse muestreando la imagen con prompts geométricos; ni los
prompts de texto ni la detección por concepto contribuyen a ese muestreo exhaustivo. Su valor
sería otro: si SAM3 resultara válido como segmentador automático, su detección y su seguimiento
integrados abrirían la puerta a unificar bajo un solo modelo los tres pasos que OVO hoy resuelve
por separado. A cambio, su backbone, sensiblemente mayor que el de SAM2, incrementa el coste
computacional y la memoria, magnitudes que se cuantifican más adelante (Figura~\ref{fig:seg_cost}). La
Tabla~\ref{tab:seg_sam2_vs_sam3} resume estas diferencias; ambas hipótesis, mejores máscaras y
una posible unificación, se examinan a continuación, empezando por la más básica: la cantidad y
la calidad de las máscaras.

\begin{table}[t]
  \centering
  \caption{Main differences between SAM2 and SAM3 as segmenters.}
  \label{tab:seg_sam2_vs_sam3}
  \begin{tabular}{lll}
    \toprule
     & \textbf{SAM2} & \textbf{SAM3} \\
    \midrule
    Vision backbone         & Hiera-L              & Perception Encoder (PE) \\
    Backbone parameters     & $\sim$214M           & $\sim$816M \\
    Pre-training            & Masked autoencoding  & Image-text contrastive \\
    Language-aligned        & No                   & Yes \\
    Text prompts            & No                   & Yes \\
    Automatic detection     & No                   & Yes (DETR + presence head) \\
    Video tracking          & Yes                  & Yes (inherited from SAM2) \\
    Automatic generation (AMG) & Yes               & No (built here, Sec.~\ref{sec:seg_results}) \\
    \bottomrule
  \end{tabular}
\end{table}

\subsection{Effect and cost of a newer segmenter}
\label{sec:seg_results}

A diferencia de SAM2, SAM3 no proporciona un generador automático de máscaras. Para evaluarlo como segmentador de OVO sin alterar
el resto del sistema, se ha construido un AMG para SAM3. La estrategia preserva la exigencia de
no cambiar OVO: se reutiliza sin modificar la maquinaria del AMG de SAM2, es decir, la
generación de la rejilla, el procesado por lotes y la poda, y se sustituye únicamente el
predictor por el interactivo de SAM3. De este modo, el preprocesado de imagen (redimensionado a
1008 y normalización) y los módulos que producen las máscaras (prompt encoder y mask decoder)
son los de SAM3, mientras que la orquestación exterior, puramente geométrica y agnóstica al
modelo, es la de SAM2.

La evaluación se plantea en dos frentes. El primero opera a nivel de punto y examina la
ambigüedad de granularidad descrita en la Sección~\ref{sec:seg_amg}. La
Figura~\ref{fig:point_ambiguity} pincha un mismo punto, sobre la mesa contigua a la pantalla del
frame 70, en ambos modelos. SAM2 responde con una máscara de confianza alta (0.969) que fusiona
pantalla y mesa, el mismo error que se veía en la Figura~\ref{fig:amg_pipeline}; SAM3 evita esa
fusión, pero en lugar de recuperar la mesa entera segmenta solo una parte, pasándose al extremo
opuesto. Ninguno resuelve la ambigüedad: un clic no determina la granularidad pretendida, y el
modelo más reciente la reparte de otra forma, pero no la elimina.

\begin{figure}[htbp]
  \centering
  \includegraphics[width=\textwidth]{point_ambiguity}
  \caption[Ambigüedad de granularidad a nivel de punto]{Un mismo punto pinchado en SAM2 y SAM3
    sobre el frame 70 de \emph{office0}. Cada modelo devuelve tres máscaras multimask con su
    confianza; ninguna versión desambigua la granularidad. SAM2 fusiona pantalla y mesa, SAM3
    solo recorta una parte.}
  \label{fig:point_ambiguity}
\end{figure}

El segundo frente opera a nivel de AMG, es decir, SAM3 actuando como segmentador por frame de
OVO. En cuanto a cantidad, sobre el frame representativo SAM3 produce más máscaras que SAM2,
30 en bruto frente a 23, y 23 frente a 19 tras la poda de OVO. La Figura~\ref{fig:seg_scene}
confirma que la diferencia se mantiene a lo largo de \emph{office0}, con una media de 21.8
máscaras finales frente a 20.4. Ese recuento, sin embargo, engaña: la cobertura, la fracción de
imagen que las máscaras finales cubren, es sistemáticamente menor con SAM3, un 51.4\% de media
frente al 60.7\% de SAM2. Es decir, SAM3 genera más máscaras pero cubre menos escena, la
señal más clara de que las máscaras de más son fragmentos y no objetos nuevos. En cuanto a
calidad, las segmentaciones finales difieren de forma
visible, como ilustra la Figura~\ref{fig:amg_sam3}. SAM3 tiende sistemáticamente a una
granularidad mayor: donde SAM2 entrega una máscara por objeto, SAM3 lo trocea en varias, como el
sofá del frame 750. Esa granularidad fina corrige algunos de los defectos de SAM2, por ejemplo
en el frame 1930, donde SAM2 fusiona dos pantallas en una sola máscara y SAM3 las separa; pero
la mayoría de las veces sobrepasa el objeto y lo fragmenta, de modo que las máscaras de más no
son objetos nuevos, sino partes de uno ya presente.

\begin{figure}[htbp]
  \centering
  \includegraphics[width=\textwidth]{amg_sam3}
  \caption[AMG con SAM2 frente a SAM3]{Generación automática de máscaras con SAM2 y con SAM3
    sobre tres frames de \emph{office0}. SAM3 produce de forma sistemática una granularidad más
    fina: trocea en varias máscaras lo que SAM2 entrega como una (frames 750 y 1310) y, aunque a
    veces separa objetos que SAM2 fusiona (las dos pantallas del frame 1930), en general
    sobresegmenta. Los casos numerados se comentan en el texto.}
  \label{fig:amg_sam3}
\end{figure}

\begin{figure}[htbp]
  \centering
  \includegraphics[width=\textwidth]{seg_scene}
  \caption[Máscaras, cobertura y tiempo a lo largo de office0]{Número de máscaras finales, cobertura
    de la imagen y tiempo por frame a lo largo de \emph{office0}, muestreando cada 10 frames. SAM3
    mantiene un recuento de máscaras algo mayor que SAM2 pero una cobertura consistentemente menor
    y un coste por frame mayor: más máscaras, sobre menos superficie y más caras.}
  \label{fig:seg_scene}
\end{figure}

El paso definitivo es medir el efecto sobre el mapa. Se ejecuta OVO completo con SAM3 como único
cambio, manteniendo Vanilla, CLIP y los umbrales idénticos a la baseline de SAM2, sobre las ocho
escenas de Replica, y se comparan las métricas de instancia. La Tabla~\ref{tab:seg_end2end}
recoge el resultado: SAM3 no mejora ninguna de las tres métricas; de media empeora el mIoU en
3.7~puntos, la mAcc en 2.7 y el AP agnóstico en 3.0. La sobre-segmentación observada en 2D se
traslada al mapa: las máscaras de más fragmentan instancias en lugar de descubrir objetos, y esa
fragmentación degrada tanto la etiqueta semántica como la delimitación de las instancias.

Y el cambio tampoco sale gratis. Una pasada de AMG con SAM3 es del orden de 1.4 veces más lenta
que con SAM2 y duplica la memoria de vídeo en pico (Figura~\ref{fig:seg_cost}); el sobrecoste se
concentra por completo en el encoder, cuyo tiempo se duplica, mientras que el filtrado, ajeno al
modelo, no varía. Es un coste que se paga en balde: el peso de SAM3 viene de sus capacidades
lingüísticas y de detección, y ninguna interviene en el barrido geométrico por rejilla, de modo
que la mitad del modelo permanece inactiva mientras segmenta.

\begin{figure}[htbp]
  \centering
  \includegraphics[width=\textwidth]{seg_cost}
  \caption[Coste de una pasada de AMG: SAM2 frente a SAM3]{Coste por frame del AMG sobre el frame
    70 de \emph{office0}, aislando cada modelo. Izquierda: desglose del tiempo en encoder, decoder
    y filtrado. Derecha: memoria de vídeo en pico. SAM3 es más lento, sobre todo en el encoder, y
    ocupa el doble de memoria, casi toda peso de modelo.}
  \label{fig:seg_cost}
\end{figure}

\begin{table}[htbp]
  \centering
  \caption[SAM2 frente a SAM3 de extremo a extremo]{Métricas de mapa de OVO con SAM2 y con SAM3
    como segmentador, media sobre las ocho escenas de Replica (Vanilla + CLIP; el único cambio es
    el segmentador). Valores en~\%. AP\textsubscript{agn} es el AP de instancia class-agnostic.}
  \label{tab:seg_end2end}
  \begin{tabular}{lccc}
    \toprule
     & \textbf{mIoU} & \textbf{mAcc} & \textbf{AP\textsubscript{agn}} \\
    \midrule
    SAM2 (baseline) & 29.9 & 39.7 & 23.3 \\
    SAM3            & 26.2 & 37.0 & 20.3 \\
    \midrule
    $\Delta$        & $-3.7$ & $-2.7$ & $-3.0$ \\
    \bottomrule
  \end{tabular}
\end{table}

De aquí se sigue la conclusión del capítulo. Sustituir el modelo del segmentador altera qué
máscaras sobreviven y a qué coste, pero no elimina el techo identificado al inicio de este
capítulo. Ese techo es una propiedad del propio paradigma de generación
automática, una rejilla de prompts de un punto, ambiguos e individuales, seguida de una poda
heurística, y es externo al modelo: reside en la maquinaria de generación y poda, no en el
encoder ni en el decoder. Un modelo más reciente y más grande baraja de otra manera qué
máscaras se conservan, pero la máscara sigue siendo una apuesta de un solo intento que la
asignación hereda sin revisar. Escapar de ese techo exige actuar donde la decisión se toma
realmente, a nivel de punto, en 3D, acumulando evidencia a lo largo del tiempo en lugar de
confiar en una única máscara 2D. Ese es el objeto del capítulo siguiente.


\section{From instance to point}
\label{sec:contest_granularity}
% Por qué un fallo PARCIAL exige una decisión PARCIAL. La decisión a nivel de instancia
% es todo-o-nada, y por eso no alcanza los fallos de \ref{sec:fm_scope_limit}.
% El punto es la unidad de verdad.

Como ya se ha visto, OVO separa dos tareas: la reconstrucción \emph{geométrica} del
mapa, que corre el SLAM, y una capa \emph{semántica} de asignación y clasificación que se
monta encima de esa reconstrucción. Cada \foreignlanguage{english}{keyframe} nuevo aporta
geometría, pero la reconstrucción no acumula todo lo que ve sin criterio: para mantener la
coherencia del mapa, antes de añadir puntos comprueba qué parte de lo observado
corresponde a superficie ya reconstruida y la descarta. Solo las regiones vistas por
primera vez generan puntos nuevos; lo ya reconstruido no se duplica, únicamente anota que
se ha vuelto a observar. El resultado es una nube \emph{deduplicada}: un punto por
superficie física, no uno por observación.

Es una decisión de diseño sana, y por dos motivos. Acota la escala del mapa al área de la
escena en vez de a la longitud de la trayectoria: una pared vista cien veces sigue siendo
una pared, no cien capas de puntos. Y evita apilar reconstrucciones ligeramente
desalineadas de la misma superficie, que la emborronarían. Esa permanencia geométrica está
bien ganada: la posición de una superficie es un \emph{hecho} que no cambia, y tratarla
como fija es lo correcto.

Sobre esa reconstrucción se monta la capa semántica. Cada punto lleva un dueño, la
instancia a la que pertenece, y ese dueño lo fija la primera máscara de SAM que reunió
votos suficientes para reclamar el punto, ya sea creando una instancia nueva o adjudicándolo
a una que ya existía y bajo la que cayó. La capa semántica hereda así la permanencia de la
geometría: una vez asignado, el punto conserva su dueño durante toda la vida del mapa. Pero las dos permanencias no valen lo mismo. La
posición de un punto es un hecho medido; su dueño es una hipótesis, la de una sola máscara
en un solo instante, que nunca se revisa.

De ahí la rigidez. Por muchos \foreignlanguage{english}{keyframes} que después muestren que
el dueño correcto es otro, no hay palanca para cambiarlo. El dedup lo agrava: la superficie
ya existe, así que la máscara que la reclamaría más tarde reproyecta sobre puntos
\emph{ajenos} y no crea puntos propios con los que competir. La disputa nunca aflora como
geometría; queda sepultada en el instante de la primera asignación.

Por eso el mecanismo de fusión de \foreignlanguage{english}{loop closure}
(\ref{sec:fm_scope_limit}) no puede alcanzar estos fallos. Ese \emph{merge} une dos
instancias cuando sus nubes se solapan, y el dedup impide que ese solape exista: los puntos
mal adjudicados no tienen copia correcta encima, son los mismos puntos con el dueño
equivocado. Sin dos nubes que comparar, el merge no tiene de qué dispararse. Tres
situaciones, recurrentes en la práctica, lo ejemplifican (Fig.~\ref{fig:contest_cases}).

\begin{figure}[htbp]
  \centering
  % PLACEHOLDER: sustituir cada \fbox por \includegraphics de un screenshot de Rerun (PNG),
  % cada uno un montaje Antes | Después sobre nube coloreada por instancia.
  % Escena + IDs de instancia por fijar (candidato para reusar como figura de fallo del Ch1).
  \newcommand{\caseph}[1]{\fbox{\parbox[c][2.2cm][c]{0.96\linewidth}{\centering\footnotesize\itshape #1}}}
  \begin{subfigure}{\linewidth}
    \centering\caseph{Antes: trozo contaminado \quad$\vert$\quad Después: trozo devuelto}
    \caption{Nacimiento bajo una máscara ajena.}
    \label{fig:contest_case_a}
  \end{subfigure}\\[4pt]
  \begin{subfigure}{\linewidth}
    \centering\caseph{Antes: revisita partida en dos instancias \quad$\vert$\quad Después: reunidas}
    \caption{Revisita parcial.}
    \label{fig:contest_case_b}
  \end{subfigure}\\[4pt]
  \begin{subfigure}{\linewidth}
    \centering\caseph{Antes: fragmento huérfano en el suelo \quad$\vert$\quad Después: reabsorbido}
    \caption{Granularidad inconsistente del segmentador.}
    \label{fig:contest_case_c}
  \end{subfigure}
  \caption{Los tres modos de fallo fuera del alcance del \emph{merge} (izquierda), resueltos
  por el contest (derecha).}
  \label{fig:contest_cases}
\end{figure}

\paragraph{Caso A: nacimiento bajo una máscara ajena.}
Un objeto entra en el campo de visión y SAM lo segmenta junto a otro que no es el suyo:
una sola máscara cubre ambos. Los puntos de los dos objetos nacen bajo la misma
instancia. Más tarde, cuando el objeto se ve despejado y SAM lo segmenta bien, la nueva
instancia se crea correctamente, pero el trozo que quedó atrapado en la primera ya tiene
dueño y no se lo devuelve nadie. La contaminación es permanente (Fig.~\ref{fig:contest_case_a}).

\paragraph{Caso B: la revisita parcial.}
Un objeto se ve parcialmente y desaparece de la escena antes de consolidarse. Al reentrar
por otro lado, la vista es mejor y se le asigna una instancia nueva que crece sana. Aunque
en algún \foreignlanguage{english}{keyframe} posterior las máscaras de SAM segmenten bien
el trozo que ya se había visto la primera vez, esos puntos siguen perteneciendo a la
instancia antigua. Las dos deberían ser una, pero no se solapan (son la misma superficie
vista en dos momentos, no dos copias en el mismo sitio), así que el merge no las
considera candidatas y jamás las une (Fig.~\ref{fig:contest_case_b}).

\paragraph{Caso C: granularidad inconsistente del segmentador.}
SAM no segmenta con una granularidad estable: a veces hila muy fino y parte una superficie
grande (el suelo es el caso típico), otras la da entera. Si el primer avistamiento de una
zona la parte en exceso y el resto de la trayectoria la ve completa, el fragmento inicial
queda como una instancia diminuta \emph{dentro} de la grande. Cae dentro de ella en todos
los frames siguientes, pero al tener ya dueño no se reabsorbe: el merge, de nuevo, no ve
duplicados que unir (Fig.~\ref{fig:contest_case_c}).

Los tres comparten raíz, y no es una heurística mal calibrada: son fallos \emph{fuera del
dominio} de los mecanismos existentes. Cada uno cumple su responsabilidad, el dedup la
coherencia geométrica y la fusión de \foreignlanguage{english}{loop closure} las
duplicidades por deriva o revisita, pero ninguno \emph{refina} las asignaciones a la luz de
lo que se observa después. Ese es el hueco que llena el contest, y no pasa por deshacer el
dedup sino por hacer revisable el dueño de un punto por evidencia en vez de fijo por la
primera máscara: la geometría sigue congelada, la semántica deja de estarlo. El
refinamiento descansa sobre dos decisiones.

\begin{enumerate}
  \item \textbf{Granularidad de la decisión.}
  Un fallo parcial solo se arregla con una decisión parcial. Si la unidad de decisión es la
  instancia entera, «este trozo es de B» no se puede ni plantear: o se mueve toda la instancia
  o ninguna. Bajando la decisión al punto, cada punto puede cambiar de dueño por su cuenta, y
  el \emph{merge} clásico pasa a ser el caso extremo en que resulta que todos los puntos de
  una instancia migran a otra.

  \item \textbf{Criterio para revocar.}
  Precisamente porque la primera decisión puede estar mal, revocarla a la ligera, con la
  evidencia de un solo frame, reintroduce el mismo problema en sentido contrario. La única
  base fiable para mover un punto de dueño es la \emph{evidencia acumulada}: observar a lo
  largo de muchos \foreignlanguage{english}{keyframes} qué instancia reclama ese punto de
  forma consistente, y solo entonces decidir. De confiar ciegamente en cada máscara se pasa a
  confiar en el historial de máscaras.
\end{enumerate}

\medskip
\noindent\emph{Precondición.} Todo esto presupone geometría estable y
\foreignlanguage{english}{tracking} fiable. La evidencia se acumula reproyectando máscaras
sobre los puntos, así que si el mapa deriva o el \foreignlanguage{english}{tracking} falla,
esa reproyección cae donde no toca y genera disputas que no responden a ningún conflicto
real: la historia de un punto deja de ser comparable consigo misma y el razonamiento se
rompe. Pero esto es, de nuevo, una responsabilidad ajena. Mantener la geometría sana bajo
los puntos es tarea del \foreignlanguage{english}{refresh} geométrico y los
\foreignlanguage{english}{submaps} (\ref{ch:geometric_consistency}); el contest da esa
condición por resuelta y se ocupa solo de lo suyo, arbitrar la disputa.

\section{The contest mechanism}
\label{sec:contest_mechanism}

El contest es el mecanismo que asume la responsabilidad que faltaba: registra los conflictos
entre instancias, acumula la evidencia de cada uno y, a partir de esa evidencia, decide qué
hacer. Todo él se organiza en torno a un par de roles. Cuando dos instancias se disputan un
punto, la que lo posee actualmente es el \emph{defender} $D$ y la que se lo reclama es el
\emph{challenger} $C$. La unidad sobre la que el mecanismo razona y decide no es una
instancia suelta, sino el par ordenado $(D, C)$: toda la maquinaria (la evidencia que
almacena, las señales que calcula, la decisión que emite) está referida a ese par. Su salida
es un \emph{veredicto} por par, $v(D, C) \in \{\textsc{merge}, \textsc{split},
\textsc{no-action}\}$, que traduce la evidencia acumulada en una acción concreta sobre el
mapa o en la ausencia de ella.

El mecanismo se descompone en tres etapas, cada una responsable de una cosa y ninguna más.
El \emph{registro} (\ref{sec:contest_registry}) anota la disputa punto a punto y la agrega
por par, sin decidir. El \emph{discriminador} (\ref{sec:contest_discriminator}) lee esa
evidencia y emite el veredicto, sin tocar el mapa. El \emph{actuador}
(\ref{sec:contest_actuator}) ejecuta el veredicto sobre el mapa. Registro y discriminador
son puros, y esa pureza es deliberada: la disputa se contabiliza y se juzga por separado de
donde se aplica, y solo la última etapa modifica nada (Fig.~\ref{fig:contest_pipeline}).

\begin{figure}[ht]
  \centering
  % PLACEHOLDER: diagrama del pipeline del mecanismo, a dibujar (TikZ o drawio -> PDF).
  % Tres cajas en cadena, con inputs/outputs explícitos entre ellas:
  %   [pasos de tracking: máscaras SAM reproyectadas]
  %      -> REGISTRO -> (señales por par: containment, confidence, exclusivity, ...)
  %      -> DISCRIMINADOR -> (veredicto v(D,C) in {merge, split, no-action})
  %      -> ACTUADOR -> (mapa mutado: instancias fusionadas / puntos reasignados)
  % Marcar registro y discriminador como PUROS (no tocan el mapa); solo el actuador muta.
  \fbox{\parbox[c][3.2cm][c]{0.9\linewidth}{\centering\footnotesize\itshape
  Pipeline del contest: tracking $\to$ registro (señales por par) $\to$ discriminador
  (veredicto) $\to$ actuador (mapa mutado). Registro y discriminador puros; solo el actuador
  toca el mapa.}}
  \caption{Las tres etapas del contest y el flujo de datos entre ellas.}
  \label{fig:contest_pipeline}
\end{figure}

\subsection{The registry}
\label{sec:contest_registry}

El registro es \emph{online}: se construye a la vez que el mapa, sin pasada posterior. En
cada paso de \foreignlanguage{english}{tracking}, cada máscara de SAM se reproyecta sobre los
puntos visibles y resuelve a una instancia, la mayoritaria entre los que ya tienen dueño bajo
la máscara. Sea $p$ un punto de dueño $D = o(p)$; si cae bajo una máscara que resuelve a
$C \neq D$, hay un \emph{robo} (\foreignlanguage{english}{grab}) de $p$ por $C$ contra $D$
(Fig.~\ref{fig:contest_grab}). No se aplica: se anota, sumando un robo a la cuenta de $C$
sobre ese punto.

\begin{figure}[ht]
  \centering
  % PLACEHOLDER: diagrama esquemático del robo, a dibujar (TikZ o drawio -> PDF).
  % Contenido: cámara y máscara SAM reproyectadas sobre la nube de puntos; un punto cuyo
  % dueño actual (defender) difiere de la instancia a la que resuelve la máscara (challenger);
  % flecha que marca el grab que se anota. Opcional: partición leal / robo / huérfano.
  \fbox{\parbox[c][2cm][c]{0.85\linewidth}{\centering\footnotesize\itshape
  Diagrama del robo: máscara SAM reproyectada sobre la nube; un punto cuyo defender difiere
  del challenger de la máscara; se anota el grab.}}
  \caption{El robo: la observación atómica del contest.}
  \label{fig:contest_grab}
\end{figure}

El registro tiene dos niveles. El primero es \emph{por punto}: tres contadores indexados por
el id permanente de cada punto (Tabla~\ref{tab:contest_primary}), la verdad de base que gotea
paso a paso. Cada punto lleva su cuenta de robos (por challenger), de reclamos y de
avistamientos, y nada más se guarda a este nivel; todo lo demás se deriva de aquí.

\begin{table}[ht]
  \centering
  \small
  \begin{threeparttable}
  \renewcommand{\arraystretch}{1.5}
  \begin{tabular}{@{}l l p{7.4cm}@{}}
    \toprule
    \textbf{Signal} & \textbf{Store} & \textbf{Steps counted} \\
    \midrule
    \textbf{\texttt{grabs}}     & $\texttt{point\_id}\mapsto\{C\!:\#\mathrm{KF}\tnote{*}\}$ &
      The point, owned by $D$, fell under $C$'s mask (a grab) \\
    \textbf{\texttt{claims}}    & $\texttt{point\_id}\mapsto\#\mathrm{KF}$ &
      The point fell under any used mask ($C$ or $D$) \\
    \textbf{\texttt{sightings}} & $\texttt{point\_id}\mapsto\#\mathrm{KF}$ &
      The point reprojected into view, with or without a mask \\
    \bottomrule
  \end{tabular}
  \begin{tablenotes}
    \footnotesize\item[*]\itshape $\#\mathrm{KF}$: number of
      \foreignlanguage{english}{keyframe} steps.
  \end{tablenotes}
  \end{threeparttable}
  \caption{Primary signals: per-point counters accumulated in the store.}
  \label{tab:contest_primary}
\end{table}

El segundo nivel agrega esos contadores \emph{por par} $(D, C)$, la unidad de la decisión.
Pero antes de agregar hay que separar la señal del ruido, y de eso se ocupa la
\textbf{firmeza} (\foreignlanguage{english}{firmness}): un robo de un punto cuenta como firme
solo si el challenger se lo llevó en suficientes pasos (\emph{evidencia}, un mínimo de robos)
y esos robos son una fracción real de los reclamos del punto (\emph{compromiso}, la razón
robos/reclamos por encima de un umbral). Un roce de refilón, o un parpadeo de pocos pasos, no
es firme. El número de puntos firmes de un par es \texttt{firm\_points}, y sobre ese
subconjunto se construyen todas las señales secundarias (Tabla~\ref{tab:contest_secondary}).
Estas viven cacheadas por defender: un robo solo ensucia a su defender, cuyas señales se
recomputan en el siguiente refresco, sin re-derivar el resto.

\begin{table}[ht]
  \centering
  \small
  \begin{threeparttable}
  \renewcommand{\arraystretch}{1.5}
  \begin{tabular}{@{}l p{11.6cm}@{}}
    \toprule
    \textbf{Signal} & \textbf{How it is computed} \\
    \midrule
    \textbf{\texttt{firm\_points}} & points of $D$ that $C$ steals firmly \\
    \textbf{\texttt{total\_grabs}} & sum of all $C$'s grabs over $D$'s points (firm or not) \\
    \textbf{\texttt{containment}}  & $\mathrm{firm\_points}\,/\,|D|\tnote{*}$ \\
    \textbf{\texttt{confidence}}   & per-point $\mathrm{grabs}_C/\mathrm{claims}$, averaged over the disputed points \\
    \textbf{\texttt{exclusivity}}  & firm points disputed only by $C$ $\,/\,\mathrm{firm\_points}$ \\
    \bottomrule
  \end{tabular}
  \begin{tablenotes}
    \footnotesize\item[*]\itshape $|D|$: number of points of instance $D$.
  \end{tablenotes}
  \end{threeparttable}
  \caption{Secondary signals: per-pair $(D,C)$ aggregates over firm points.}
  \label{tab:contest_secondary}
\end{table}

Este es todo el registro del contest: un flujo de robos por punto que gotea en cada paso,
tres contadores primarios por punto y un agregado por par vivo por defender. En ningún
momento se ha tocado el mapa ni se ha revocado una asignación; el registro es una
contabilidad fiel de la disputa, deliberadamente separada de la decisión. Convertir esa
evidencia en un veredicto es competencia del discriminador.

\subsection{The discriminator}
\label{sec:contest_discriminator}

El discriminador toma los pares $(D, C)$ de un defender, con sus señales ya calculadas, y
emite el veredicto. Es una función pura: no mira el mapa ni el descriptor, solo la evidencia
que el registro acumuló. Toda ella se ordena en torno a una única cantidad, la contención
$\mathrm{containment}(D, C)$: qué fracción del defender le roba el challenger con firmeza. La
contención decide la \emph{forma} de la acción, y la \foreignlanguage{english}{confidence} y
la exclusividad deciden
si hay \emph{evidencia suficiente} para ejecutarla.

Antes de razonar nada, un filtro de masa descarta el ruido: un par con menos de
$\tau_{\text{masa}}$ robos crudos ($\mathrm{total\_grabs} < 50$) no es una disputa, es un
roce, y se ignora. Sobre los pares que sobreviven, la contención reparte la decisión en tres
bandas.

\begin{enumerate}
  \item \textbf{Contención alta} ($\mathrm{containment} \geq 0.6$): el challenger se lleva casi
  todo el defender, luego el defender es un \emph{fragmento} suyo (casos A y C). La acción
  natural es \textsc{merge}, pero solo si hay un dueño fiable: se exige una
  \foreignlanguage{english}{confidence} $\mathrm{confidence} \geq 0.5$ (de media, cada punto se roba
  en la mayoría de sus reclamos) y que
  todos los challengers firmes sean, vía la fusión, una sola instancia. Si lo contienen
  varios objetos distintos no hay dueño claro y se emite \textsc{no-action}.

  \item \textbf{Contención parcial} ($0.4 \leq \mathrm{containment} < 0.6$): el challenger se lleva
  una parte grande, pero no tanta como para que la sola magnitud baste. Aquí el umbral sube:
  \textsc{merge} solo si el trozo es \emph{limpio}, es decir con confidence alta
  ($\mathrm{confidence} \geq 0.6$) y exclusivo ($\mathrm{exclusivity} \geq 0.8$, casi todos sus puntos los
  disputa este challenger y nadie más), y de nuevo con raíz única. La evidencia de magnitud
  es más débil, así que se compensa exigiendo más limpieza en las otras dos señales.

  \item \textbf{Contención baja} ($\mathrm{containment} < 0.4$): el defender \emph{retiene} la
  mayoría de sus puntos, pero cede trozos exclusivos a uno o varios vecinos. No es un
  fragmento de nadie, es un objeto que arrastra pedazos de otros (caso B). La acción es
  \textsc{split}, y no una sino tantas como haga falta: por cada challenger con un trozo
  con confidence alta ($\mathrm{confidence} \geq 0.6$) y muy exclusivo ($\mathrm{exclusivity} \geq 0.9$) se
  desprende ese trozo hacia él. Si ningún challenger cualifica, el par está en el borde y se
  emite \textsc{no-action}.
\end{enumerate}

El umbral de exclusividad sube de $0.8$ en el \textsc{merge} parcial a $0.9$ en el
\textsc{split}, y no por capricho: un split parte geometría en dos y es más caro de deshacer
si sobra, así que se exige que la frontera del trozo sea más nítida antes de cortar. El suelo
$\mathrm{containment} = 0.4$ que separa split de merge es empírico, fijado sobre el conjunto
de \foreignlanguage{english}{tuning} (\ref{sec:contest_results}).

Visto en conjunto, el \textsc{merge} clásico reaparece como el extremo de esta escala: el
caso en que la contención es total y el defender migra entero. El discriminador solo extiende
esa misma idea a las contenciones intermedias y bajas, donde la decisión deja de ser mover
toda la instancia y pasa a mover exactamente los puntos que la evidencia respalda.

\subsection{The actuator}
\label{sec:contest_actuator}
% Ejecuta el veredicto sobre el mapa: MERGE (union-find de instancias) y SPLIT (mueve los
% split_points al challenger). Última etapa, la única que muta. Aquí se cuela el timing/coste
% real (Real-time ya no es sección propia): coste por refresco, cuello de botella y rediseño
% online. Doc: docs/contest_timing_analysis.html

\section{Results}
\label{sec:contest_results}

Los umbrales del discriminador son empíricos, así que todos los números se reportan sobre un
conjunto de escenas que el ajuste nunca vio. Las ocho escenas de Replica se parten en
\foreignlanguage{english}{tuning} (office0, office3, office4, room2), las únicas sobre las que
se fijan los umbrales, y \foreignlanguage{english}{held-out} (office1, office2, room0, room1),
que se mantienen cerradas durante todo el ajuste y se abren una sola vez, al final. Ningún
umbral se toca mirando el \foreignlanguage{english}{held-out}; el $\mathrm{AP}$ que se defiende
es siempre el suyo, de modo que mide generalización y no amoldamiento.

% PRELIMINAR: números a recalcular antes de la entrega (placeholder para ver la maqueta).
Sobre la trayectoria limpia de Replica, sin saltos inyectados, el contest sube el
$\mathrm{AP}_{\text{agnostic}}$ global de $0.234$ a $0.249$ ($+0.015$, $+6.4\%$) respecto al
mapa crudo (\emph{baseline} sin fusión sobre GT simulado). La mejora aparece en las ocho
escenas, y se sostiene en el \foreignlanguage{english}{held-out}: la ganancia media es de
$+9.7\%$ en las escenas de \foreignlanguage{english}{tuning} y de $+4.55\%$ en las de
\foreignlanguage{english}{held-out}, positiva en ambos casos
(Tabla~\ref{tab:contest_ap_scenes}).

\begin{table}[ht]
  \centering
  \small
  \renewcommand{\arraystretch}{1.2}
  \begin{tabular}{@{}l r r r r@{}}
    \toprule
    \textbf{Scene} & \textbf{Baseline} & \textbf{Contest} & \textbf{$\Delta$ abs} &
      \textbf{$\Delta$ \%} \\
    \midrule
    \multicolumn{5}{@{}l}{\textbf{Tuning}} \\
    office0 & 0.199 & 0.220 & $+0.021$ & $+10.6\%$ \\
    office3 & 0.204 & 0.210 & $+0.006$ & $+2.9\%$  \\
    office4 & 0.236 & 0.280 & $+0.044$ & $+18.6\%$ \\
    room2   & 0.232 & 0.247 & $+0.015$ & $+6.5\%$  \\
    \addlinespace[1pt]
    \textbf{mean} & \textbf{0.218} & \textbf{0.239} & $\mathbf{+0.021}$ & $\mathbf{+9.7\%}$ \\
    \midrule
    \multicolumn{5}{@{}l}{\textbf{Held-out}} \\
    office1 & 0.170 & 0.175 & $+0.005$ & $+2.9\%$  \\
    office2 & 0.210 & 0.225 & $+0.015$ & $+7.1\%$  \\
    room0   & 0.269 & 0.287 & $+0.018$ & $+6.7\%$  \\
    room1   & 0.344 & 0.349 & $+0.005$ & $+1.5\%$  \\
    \addlinespace[1pt]
    \textbf{mean} & \textbf{0.248} & \textbf{0.259} & $\mathbf{+0.011}$ & $\mathbf{+4.55\%}$ \\
    \midrule
    \textbf{All} & \textbf{0.234} & \textbf{0.249} & $\mathbf{+0.015}$ & $\mathbf{+6.4\%}$ \\
    \bottomrule
  \end{tabular}
  \caption{$\mathrm{AP}_{\text{agnostic}}$ per Replica scene, baseline vs.\ contest.}
  \label{tab:contest_ap_scenes}
\end{table}

Las métricas \emph{class-aware} (mIoU / mAcc) no empeoran: se mueven ligeramente al alza
(mIoU $+0.2$, mAcc $+0.3$ globales, Tabla~\ref{tab:contest_semantic}). Es coherente con la
naturaleza del mecanismo: el contest reorganiza la propiedad de los puntos entre instancias,
no cambia la etiqueta semántica que cada una recibe, así que su efecto vive en la granularidad
de las instancias (lo que mide el $\mathrm{AP}$ agnóstico) y deja el mIoU por clase
prácticamente intacto.

\begin{table}[ht]
  \centering
  \small
  \renewcommand{\arraystretch}{1.2}
  \begin{tabular}{@{}l r r r r r r@{}}
    \toprule
    & \multicolumn{3}{c}{\textbf{mIoU}} & \multicolumn{3}{c}{\textbf{mAcc}} \\
    \cmidrule(lr){2-4} \cmidrule(lr){5-7}
    \textbf{Scene} & Base & Contest & $\Delta$ & Base & Contest & $\Delta$ \\
    \midrule
    office0 & 25.4 & 25.9 & $+0.5$ & 33.9 & 34.3 & $+0.4$ \\
    office1 & 16.2 & 16.1 & $-0.1$ & 31.7 & 31.7 & $+0.0$ \\
    office2 & 24.6 & 25.9 & $+1.3$ & 32.4 & 34.0 & $+1.5$ \\
    office3 & 25.8 & 25.9 & $+0.0$ & 32.3 & 32.3 & $+0.0$ \\
    office4 & 43.1 & 42.7 & $-0.3$ & 55.4 & 55.8 & $+0.4$ \\
    room0   & 33.2 & 33.6 & $+0.4$ & 39.9 & 40.3 & $+0.4$ \\
    room1   & 38.6 & 40.2 & $+1.6$ & 53.8 & 56.1 & $+2.3$ \\
    room2   & 26.6 & 26.6 & $+0.0$ & 35.1 & 35.1 & $+0.0$ \\
    \midrule
    \textbf{All} & \textbf{27.1} & \textbf{27.3} & $\mathbf{+0.2}$ & \textbf{40.2} & \textbf{40.5} & $\mathbf{+0.3}$ \\
    \bottomrule
  \end{tabular}
  \caption{Class-aware semantic metrics (\%) per Replica scene, baseline vs.\ contest.}
  \label{tab:contest_semantic}
\end{table}
% Validación en AISLAMIENTO, sobre trayectoria sana y sin saltos: se atribuye la mejora
% al mecanismo, no a su interacción con el ruido. COROLARIO, que hay que escribir: el
% contest PRESUPONE geometría sana y tracking fiable -> el sustrato no es un preámbulo,
% es una PRECONDICIÓN.
%
% ABLACIÓN 2x2 (el experimento estrella de la memoria), sobre saltos inyectados:
%
%                    | sin submapas            | con submapas
%   .................+.........................+...........................
%   sin contest      | baseline degradado      | ¿el submapa solo, ayuda?
%   con contest      | contest sobre sustrato  | la propuesta completa
%                    | roto -> debería ir mal  |
%
% La celda de abajo-izquierda es la que DEMUESTRA que el orden causal es real y no
% retórica. Asegurar este experimento sí o sí.
